% Part: computability
% Chapter: recursive-functions
% Section: halting-problem

\documentclass[../../../include/open-logic-section]{subfiles}

\begin{document}

\olfileid{cmp}{rec}{hlt}
\olsection{停机问题}

广义的\emph{停机问题}是指，给定可计算函数的描述~$e$（如程序）与数~$n$，判定该函数在输入~$n$ 上的计算是否\emph{停机}（即是否产生结果）。众所周知，Alan Turing 证明了该问题本身无法由可计算函数解决，即函数
\[
h(e, n) =
\begin{cases}
1 & \text{若计算 $e$ 在输入 $n$ 上停机}\\
0 & \text{否则，}
\end{cases}
\]
是不可计算的。

在部分递归函数的语境下，程序的描述可由 Kleene（克林）范式定理中给出的指标~$e$ 充当。若 $f$ 是部分递归函数，则使得范式定理中方程成立的任意~$e$，都是~$f$ 的指标。给定数~$e$，Kleene 范式定理指出，\[
\cfind{e}(x) \simeq U(\mu s \; T(e, x, s))
\]
是部分递归的，且对任意部分递归函数 $f\colon \Nat \to \Nat$，存在 $e \in \Nat$，使得对所有~$x \in \Nat$，$\cfind{e}(x) \simeq
f(x)$。事实上，对每个这样的 $f$，这样的~$e$ 不止一个，而是有无穷多个。\emph{停机函数}~$h$ 定义如下：
\[
h(e, x) =
\begin{cases}
1 & \text{若 $\cfind{e}(x) \fdefined$}\\
0 & \text{否则。}
\end{cases}
\]
注意，当 $\cfind{e}(x) \fundefined$ 时 $h(e, x) = 0$，但当~$e$ 根本不是某个部分递归函数的指标时，$h(e, x) = 0$ 同样成立。

\begin{thm}
\ollabel{thm:halting-problem}
停机函数~$h$ 不是部分递归的。
\end{thm}

\begin{proof}
若 $h$ 是部分递归的，则我们可以定义
\[
d(y) =
\begin{cases}
1 & \text{若 $h(y, y) = 0$}\\
\umin{x}{x \neq x} & \text{否则。}
\end{cases}
\]
由于没有数 $x$ 满足 $x \neq x$，因此 $\umin{x}{x \neq x}$ 无定义，故 $d(y) \fundefined$ 当且仅当 $h(y,y) \neq 0$。由此定义可得：
\begin{enumerate}
\item $d(y) \fdefined$ 当且仅当 $\cfind{y}(y) \fundefined$ 或 $y$ 不是某个部分递归函数的指标。
\item $d(y) \fundefined$ 当且仅当 $\cfind{y}(y) \fdefined$。
\end{enumerate}
若 $h$ 是部分递归的，则 $d$ 也将是部分递归的。因此，根据 Kleene 范式定理，$d$ 有一个指标~$e_d$。
考虑 $h(e_d, e_d)$ 的值。有两种可能的情况：$0$ 和~$1$。
\begin{enumerate}
\item 若 $h(e_d, e_d) = 1$，则 $\cfind{e_d}(e_d) \fdefined$。但 $\cfind{e_d} \simeq d$，且 $d(e_d)$ 有定义当且仅当 $h(e_d, e_d) = 0$。故 $h(e_d, e_d) \neq 1$。
\item 若 $h(e_d, e_d) = 0$，则要么 $e_d$ 不是某个部分递归函数的指标，要么它是且 $\cfind{e_d}(e_d)
  \fundefined$。但同样地，$\cfind{e_d} \simeq d$，且 $d(e_d)$ 无定义当且仅当 $\cfind{e_d}(e_d) \fdefined$。
\end{enumerate}
其结果是，$e_d$ 终究不可能是一个部分递归函数的指标。然而，若 $h$ 是部分递归的，$d$ 也会是，从而将 $e_d$ 作为其指标将是合法的。由此只能得出结论：$h$ 不可能是部分递归的。
\end{proof}

\end{document}